\documentclass[12pt, twocolumn, landscape, a4paper]{article}

\usepackage[margin=1cm]{geometry}
\usepackage[utf8]{inputenc}
\usepackage[spanish]{babel}
\usepackage{amssymb, amsmath, amsbsy}
\usepackage{verbatim}
\usepackage{bookman}
\usepackage{enumerate}
% \usepackage{tikz}
% \usetikzlibrary{angles,quotes}
\usepackage[pdftex]{hyperref}
\hypersetup{ pdftitle = {   },pdfauthor = {Luis Carrillo Gutiérrez} }

\renewcommand{\familydefault}{\sfdefault}

\begin{document}

\section*{Cramer}

\noindent Sea el sistema de ecuaciones lineales con 3 incógnitas ($\,x_{i}\,$) :
\begin{eqnarray*}
a_{11}\,x_{1} + a_{12}\,x_{2} + a_{13}\,x_{3} = b_{1} \\
a_{21}\,x_{1} + a_{22}\,x_{2} + a_{23}\,x_{3} = b_{2} \\
a_{31}\,x_{1} + a_{32}\,x_{2} + a_{33}\,x_{3} = b_{3}
\end{eqnarray*}

\noindent Puede ser expresado (representación matricial o forma matricial) como una relación entre matrices ($\,\mathrm{A} \times x_{i} = b\,$):

$$ \left (
\begin{array}{ccc}
a_{11} & a_{12} & a_{13} \\
a_{21} & a_{22} & a_{23} \\ 
a_{31} & a_{32} & a_{33} 
\end{array}
\right ) \times \left ( \begin{array}{c} x_{1} \\ x_{2} \\ x_{3} \end{array} \right ) = \left ( \begin{array}{c} b_{1} \\ b_{2} \\ b_{3} \end{array} \right )
$$ \\

\noindent Se define la \emph{matriz ampliada}, \emph{completa} o \emph{aumentada} del sistema de ecuaciones lineales, como la matriz por bloques ($\,\mathrm{A}^{*} = \mathrm{A}\,|\,b\,$):

$$ \left (
\begin{array}{ccc|r}
a_{11} & a_{12} & a_{13} & b_{1} \\
a_{21} & a_{22} & a_{23} & b_{2} \\ 
a_{31} & a_{32} & a_{33} & b_{3}
\end{array}
\right )
$$ \\

\noindent Sistemas de tres ecuaciones con tres incógnitas utilizando las fórmulas de Cramer. \\

\begin{itemize}
\item{ %\begin{eqnarray*}
$\left .
\begin{array}{rc}
x + y + z &= \quad 11 \\
2x - y + z &= \quad\;5 \\ 
3x + 2y + z &= \quad 24
\end{array} \right \} \Rightarrow \left (
\begin{array}{ccc|r}
1 & 1 & 1 & 11\\
2 & -1 & 1 & 5 \\ 
3 & 2 & 1 & 24 \\
\end{array} \right )
%\end{eqnarray*}
$
}
\end{itemize}

\begin{eqnarray*}
det(A) = |\mathrm{A}| = 
\begin{vmatrix}
\,1 & 1 & 1 \\
\,2 & -1 & 1 \\ 
\,3 & 2 & 1 \\
\end{vmatrix} = (-1 + 3 + 4) - (-3 + 2 + 2) = 5
\end{eqnarray*}

\begin{eqnarray*}
|\mathrm{A}_{x}| = 
\begin{vmatrix}
\mathtt{11} & 1 & 1 \\
\mathtt{5} & -1 & 1 \\ 
\mathtt{24} & 2 & 1 \\
\end{vmatrix} = (-11 + 24 + 10) - (-24 + 22 + 5) = 20
\end{eqnarray*}

\begin{eqnarray*}
|A_{y}| = 
\begin{vmatrix}
1 & \mathtt{11} & 1 \\
2 & \mathtt{5}  & 1 \\ 
3 & \mathtt{24} & 1 \\
\end{vmatrix} = (5 + 33 + 48) - (15 + 24 + 22) = 25
\end{eqnarray*}

\begin{eqnarray*}
|A_{z}| = 
\begin{vmatrix}
1 & 1 & \mathtt{11} \\
2 & -1 & \mathtt{5} \\ 
3 & 2 & \mathtt{24} \\
\end{vmatrix} = (-24 + 15 + 44) - (-33 + 10 + 48) = 10
\end{eqnarray*}

$x = \dfrac{|A_{x}|}{|A|} = \dfrac{20}{5} = 4$ ; $\quad\quad y = \dfrac{|A_{y}|}{|A|} = \dfrac{25}{5} = 5$ ; $\quad\quad z = \dfrac{|A_{z}|}{|A|} = \dfrac{10}{5} = 2$

% \begin{comment}\end{comment}

\begin{itemize}
\item{% \begin{eqnarray*}
$\left .
\begin{array}{cc}
2x + y + 2z & = 4 \\
-x + 3y + z & = 3 \\ 
x - y &= 1
\end{array} \right \} \Rightarrow \left (
\begin{array}{ccc|r}
2 & 1 & 2 & \mathtt{4} \\
-1 & 3 & 1 & 3 \\ 
1 & -1 & 0 & 1 \\
\end{array} \right )
% \end{eqnarray*}
$
}
\end{itemize}

\begin{itemize}
\item{$
\begin{array}{ccr}
2x + 6y + z & = & 7 \\
2y - z & = & -1 \\
5x + 7y - 4z & = & 9
\end{array}
$}
\end{itemize}
\end{document}